\documentclass[10pt,a4paper]{amsart}
\usepackage[margin=19mm]{geometry}
\usepackage{amsmath,amssymb,mathtools}
\usepackage[scr=rsfs,cal=euler]{mathalfa}
\usepackage{xfrac}
\usepackage[unicode,hidelinks,bookmarks=false]{hyperref}
\newcommand{\ie}{i.e.\ }
\newcommand{\cf}{cf.\ }
\newcommand{\resp}{respectively\ }
\newcommand{\Subsection}{\subsection}
\providecommand{\nobreakdash}{\nobreak\hbox{-}}
\setlength{\emergencystretch}{2em}
\multlinegap0pt
\usepackage{bm}
\usepackage{bbm}
\usepackage{dsfont}
\usepackage{xspace}
\usepackage{cancel}
\usepackage{mathtools}
\usepackage{amsthm}
\makeatletter
\@ifundefined{theorem}{\newtheorem{theorem}{Theorem}}{}
\@ifundefined{prop}{\newtheorem{prop}[theorem]{Proposition}}{}
\makeatother
\DeclareMathOperator{\Sat}{Sat}
\def\CC{\mathbb{C}}
\title[An algorithm for annihilator and Bernstein-Sato polynomial o]{An algorithm for annihilator and Bernstein-Sato polynomial of a rational function}
\author{Manuel Gonz\'alez-Villa, Edwin Le\'on-Cardenal, Viktor Levandovskyy, Jorge Mart\'in-Morales}
\date{Source: arXiv:2602.00280v1 (CC BY 4.0)}
\begin{document}
\maketitle

\noindent\emph{Source and licence.} An algorithm for annihilator and Bernstein-Sato polynomial of a rational function. Version arXiv:2602.00280v1 (CC BY 4.0), distributed under the Creative Commons Attribution 4.0 International licence, as recorded in the arXiv licence metadata for this exact version and shown on the abstract page https://arxiv.org/abs/2602.00280v1. Full rights notice: licenses/arxiv-2602.00280-CC-BY-4.0.txt.\par\smallskip

\noindent\textit{Editorial scope.} Counted unit: the Theorem whose source label is \texttt{thm:main-theorem} in the article (source0.tex lines 492--498, source label \texttt{thm:main-theorem}), delivered together with the article's own complete original proof of it (source0.tex lines 500--576, one \texttt{proof} environment of 77 lines). No mathematics is written, repaired or re-derived: statement and proof are the article's own text line for line. Required context retained uncounted: prop:saturation (lines 438--440). Every definition, display and label of those lines is retained unchanged. Omitted from this package and not redistributed: 1--437, 441--491, 499--499, 577--1068. Nothing inside a retained excerpt is omitted or altered: every retained body is delivered exactly as the article's lines are written, so no omission receipt of any kind accompanies this package. Citations: the retained excerpts carry no citation command at all, so no bibliography entry of the article is transcribed and no reference is redistributed. Reference adaptation: none was needed and none was made; every reference inside the delivered unit resolves inside it, so no unresolved reference or citation is produced. Printed slips kept literal: no printed word, symbol, hypothesis or number of the counted statement or of its proof was altered. Layout and class: the article's own class is \texttt{amsart} with packages including geometry, amsthm, amsmath, amssymb, mathtools, algorithm2e, inputenc, babel, hyperref, xcolor, tikz, tikz-cd; the wrapper is the lane's pinned \texttt{amsart} editorial wrapper at 10pt on A4, which restates the theorem-like environments the retained text needs and adds the article's own macro definitions that the retained text needs (\texttt{\textbackslash CC}, \texttt{\textbackslash Sat}), with extra wrapper packages (bm, bbm, dsfont, xspace, cancel, mathtools). The delivered numbering is the unit's own. Assets: no retained span contains a figure environment, a graphics-inclusion command, a PDF-page inclusion, a table or a TikZ picture; no image file is copied into this package, the package declares no asset dependency and no image file is redistributed with it. Licence: version arXiv:2602.00280v1 of this article is distributed under the Creative Commons Attribution 4.0 International licence, as recorded in the arXiv licence metadata for this exact version and shown on the abstract page https://arxiv.org/abs/2602.00280v1. A full rights notice is delivered at licenses/arxiv-2602.00280-CC-BY-4.0.txt. Characters: every retained excerpt is pure ASCII, so the delivered engine, XeLaTeX with its default Latin Modern text font, reports no missing character.
\begin{prop}\label{prop:saturation}
The ideal $I_m(s)$ (resp.~$I(s_1,s_2)$) is $\CC[s]$-saturated (resp.~$\CC[s_1,s_2]$-saturated).
\end{prop}

\begin{theorem}\label{thm:main-theorem}
Assume that the Bernstein-Sato ideal $B_{f,g}$ satisfies the $\mathcal{C}_m$-condition.
Then
\[
I_m(s) = \Sat_{\CC[s]}(I(s,-s-m)).
\]
\end{theorem}

\begin{proof}
Clearly $I(s,-s-m) \subseteq I_m(s)$. Then $\Sat_{\CC[s]}(I(s,-s-m)) \subseteq \Sat_{\CC[s]}(I_m(s)) = I_m(s)$,
since $I_m(s)$ is $\CC[s]$-saturated, see Proposition~\ref{prop:saturation}.

For the other inclusion take $P(s) \in I_m(s)$.
We need to find $\tilde{P}(s_1,s_2) \in I(s_1,s_2)$ such that $\tilde{P}(s,-s-m) = q(s) P(s)$ for some
$q(s) \in \CC[s] \setminus \{0\}$.
In order to do this we use the functional equation associated with the element $b(s_1,s_2) \in B_{f,g}$
satisfying the $\mathcal{C}_m$-condition,
\begin{equation}\label{eq:BSfg}
Q(s_1,s_2) f^{s_1+1} g^{s_2+1} = b(s_1,s_2) f^{s_1} g^{s_2},
\end{equation}
where $Q(s_1,s_2) \in D[s_1,s_2]$. Let us write
\begin{equation}\label{eq:hfgk}
P(s_1) f^{s_1} g^{s_2}
= \frac{h(s_1,s_2)}{(fg)^k} f^{s_1} g^{s_2}
\ \in \ \CC \Big[ x_1,\ldots,x_n,s_1,s_2,\frac{1}{fg} \Big] f^{s_1} g^{s_2}
\end{equation}
for some $h(s_1,s_2) \in \CC[x_1,\ldots,x_n,s_1,s_2]$ and $k \geq 0$.
Applying the substitution $s_1 = s$, $s_2=-s-m$ to \eqref{eq:hfgk}, we obtain
\[
P(s) \frac{1}{g^m} \left( \frac{f}{g} \right)^s
= \frac{h(s,-s-m)}{(fg)^k} \frac{1}{g^m} \left( \frac{f}{g} \right)^s.
\]
Since $P(s) \frac{1}{g^m} \left( \frac{f}{g} \right)^s = 0$,
we get
\begin{equation}\label{eq:hzero}
h(s,-s-m) = 0.
\end{equation}
% or equivalently
% \[
% h(s_1,s_2) = (s_2 + s_1 + m) \tilde{h}(s_1,s_2)
% \]
% for some $\tilde{h}(s_1,s_2) \in \CC[x_1,\ldots,x_n,s_1,s_2]$.

If $k = 0$, then it is enough to choose $\tilde{P}(s_1,s_2) = P(s_1) - h(s_1,s_2)$ and $q(s)=1$.
Note that $\tilde{P}(s_1,s_2) \in I(s_1,s_2)$ by~\eqref{eq:hfgk} and moreover, $\tilde{P}(s,-s-m) = q(s) P(s)$ holds by~\eqref{eq:hzero}.

% If $k=1$, then~\eqref{eq:hfgk} and~\eqref{eq:BSfg} shifted by $1$ in $s_1$ and $s_2$ give rise to
% \[
% \begin{aligned}
% b(s_1-1,s_2-1) P(s_1) f^{s_1} g^{s_2}
% &= b(s_1-1,s_2-1) \frac{h(s_1,s_2)}{fg} f^{s_1} g^{s_2} \\
% &= h(s_1,s_2) b(s_1-1,s_2-1) f^{s_1-1} g^{s_2-1} \\[4pt]
% &= h(s_1,s_2) Q(s_1-1,s_2-1) f^{s_1} g^{s_2}.
% \end{aligned}
% \]
% Choose $\tilde{P}(s_1,s_2) = b(s_1-1,s_2-1) P(s_1) - h(s_1,s_2) Q(s_1-1,s_2-1) \in I_m(s_1,s_2)$
% and $\tilde{P}(s,-s-m) = b(s-1,s-m-1) P(s)$ holds by~\eqref{eq:hzero}. Note that $b(s-1,s-m-1) \neq 0$ by assumption.

If $k\geq 1$, then we multiply \eqref{eq:hfgk} by $b(s_1-k,s_2-k)$ and obtain
\[
b(s_1-k,s_2-k) P(s_1) f^{s_1} g^{s_2} = h(s_1,s_2) b(s_1-k,s_2-k) f^{s_1-k} g^{s_2-k}.
\]
Applying the functional equation~\eqref{eq:BSfg} to the right-hand side of the previous equation, we get
\[
b(s_1-k,s_2-k) P(s_1) f^{s_1} g^{s_2} = h(s_1,s_2) Q(s_1-k,s_2-k) f^{s_1-(k-1)} g^{s_2-(k-1)}.
\]
Iterating this process $k-1$ times and using~\eqref{eq:hfgk} and~\eqref{eq:BSfg} as above, we arrive at the following functional equation
\begin{equation}\label{eq:new-function-eq}
\tilde{b}(s_1,s_2) P(s_1) f^{s_1} g^{s_2} = h(s_1,s_2) R(s_1,s_2) f^{s_1} g^{s_2},
\end{equation}
\[
\mbox{where }
\tilde{b}(s_1,s_2) = \prod_{i=1}^k b(s_1-i, s_2-i) \mbox{ and } R(s_1,s_2) = Q(s_1-k,s_2-k) \cdot \ldots \cdot Q(s_1-1,s_2-1).
\]
Finally we choose
\[
\begin{aligned}
\tilde{P}(s_1,s_2) &= \tilde{b}(s_1,s_2) P(s_1) - h(s_1,s_2) R(s_1,s_2), \\
q(s) &= \tilde{b}(s,-s-m) = \prod_{i=1}^k b(s-i, -s-m-i).
\end{aligned}
\]
Therefore by~\eqref{eq:new-function-eq} we conclude that $\tilde{P}(s_1,s_2) \in I(s_1,s_2)$,
and because of~\eqref{eq:hzero} $\tilde{P}(s,-s-m) = \tilde{b}(s,-s-m) P(s)$ follows.
Moreover, $q(s) = \tilde{b}(s,-s-m) \neq 0$ because we have chosen $b(s_1,s_2)$ satisfying the $\mathcal{C}_m$-condition.
\end{proof}

\end{document}
